\documentclass[11pt,a4paper]{article}

\usepackage[utf8]{inputenc}
\usepackage[margin=1in]{geometry}
\usepackage{amsmath,amssymb,amsfonts}
\usepackage{booktabs}
\usepackage{tabularx}
\usepackage{multirow}
\usepackage{xcolor}
\usepackage{hyperref}
\usepackage{titlesec}
\usepackage{fancyhdr}
\usepackage{listings}

\hypersetup{
    colorlinks=true,
    linkcolor=blue,
    citecolor=blue,
    urlcolor=blue
}

\pagestyle{fancy}
\fancyhf{}
\rhead{\footnotesize SWR302 -- RP6 Empirical Validation Report}
\lhead{\footnotesize Student: Nguyen Pham Xuan Nhi (SE201170)}
\rfoot{\thepage}

\title{\textbf{Pilot Validation Report: Semantic Ambiguity Resolution via Low-Fidelity Prototyping on a Student Software Project}\\
\large A Controlled Audit on 10 User Stories from the Self-Storage Management System}

\author{\textbf{Nguyen Pham Xuan Nhi} (Student ID: SE201170) \\
Department of Software Engineering, FPT University, Ho Chi Minh City, Vietnam \\
Course: SWR302 (Software Requirements Engineering) --- Topic RP6}

\date{October 2026}

\begin{document}

\maketitle

\begin{abstract}
This report provides the experimental protocol, audit matrix, and statistical verification for the preliminary pilot test described in Section~VI-B of the paper \textit{The Effect of Prototyping on Identifying Missing and Unclear Requirements}. Conducted on ten user stories from an undergraduate software engineering project (Self-Storage Facility Rental and Management System), the audit examines how pairing low-fidelity wireframes with inspection checklists resolves requirements ambiguity. Across four linguistic ambiguity categories (structural, scope, anaphoric, and vagueness), fifteen pre-prototype defects were identified across all ten statements (100\% defect rate). After aligning user stories with wireframes and visual acceptance criteria, fourteen defects were eliminated, achieving an Ambiguity Resolution Rate of $\mathbf{ARR = 93.3\%}$. A non-parametric Wilcoxon signed-rank test confirmed statistical significance ($W = 55, Z = -2.81, p = 0.005 < 0.01$). The single unresolved defect involved response latency, illustrating the non-functional requirement (NFR) limitations of static prototypes.
\end{abstract}

\vspace{0.5em}
\noindent\textbf{Keywords:} Requirements Engineering, Prototyping, Semantic Ambiguity, Pilot Validation, Wilcoxon Signed-Rank Test, Gherkin Criteria, SWR302.

\vspace{1em}
\hrule
\vspace{1em}

\section{Introduction and objectives}

While systematic literature reviews summarize established research, proposed research directions often require localized validation to confirm their practicality in educational settings. To connect theoretical findings with practical requirements engineering, this report documents a pilot validation conducted on undergraduate software project requirements.

The pilot test focuses on three specific objectives: first, applying theoretical linguistic ambiguity classifications (Berry et al., 2003; Cascon et al., 2025) directly to student user stories; second, measuring the defect reduction rate achieved when requirements are paired with wireframes; and third, testing whether static visual prototypes show the non-functional requirement deficit described in empirical literature (Rueda et al., 2020; Abad et al., 2018).

\section{Experimental protocols for proposed directions}

As outlined in Section~VI of the main paper, three experimental designs were developed to investigate prototyping interventions in agile workflows.

Direction 1 examines Enriched User Stories in agile sprints. In typical agile development, user stories follow a concise format that frequently omits interface layout details and business constraints, leading to misunderstandings and rework during demonstrations. In this proposed workflow, designers and product owners create low-fidelity wireframes one sprint ahead of development ($n-1$) and attach them directly to backlog items alongside visual acceptance criteria. Key evaluation metrics include sprint velocity, developer clarification questions logged during sprints, and total rework hours after sprint reviews.

Direction 2, which forms the basis of this pilot study, evaluates a manual ambiguity checklist paired with interface wireframes. Requirements statements are inspected against four formal linguistic categories: structural ambiguity, scope ambiguity, anaphoric ambiguity, and vagueness. Analysts and stakeholders check each statement against screen wireframes to determine whether interface elements clarify ambiguous terms. Efficacy is quantified using the Ambiguity Resolution Rate:
\begin{equation}
    \text{ARR} = \frac{D_{\text{pre}} - D_{\text{post}}}{D_{\text{pre}}} \times 100\%
\end{equation}
where $D_{\text{pre}}$ and $D_{\text{post}}$ denote the total number of ambiguity defects identified before and after prototype alignment.

Direction 3 evaluates wireframe walkthroughs augmented with a think-aloud protocol. Because passive paper mockups struggle to capture performance and security expectations, participants verbalize their thought process during interactive tasks while facilitators prompt them regarding network interruptions, data confidentiality, and latency tolerances.

\newpage
\section{Pilot audit matrix on user stories}

The pilot study analyzed ten user stories from the unrefined software requirements specification of the Self-Storage Facility Rental and Management System, an undergraduate capstone project. Table~\ref{tab:audit_matrix} documents each statement, the identified ambiguity defects, the wireframe alignment mechanism, and the resulting resolution status.

\begin{table*}[htbp]
\centering
\scriptsize
\caption{Ambiguity Audit Matrix: Pre- vs. Post-Prototype Inspection}
\label{tab:audit_matrix}
\begin{tabularx}{\textwidth}{p{0.07\textwidth} p{0.22\textwidth} p{0.24\textwidth} p{0.25\textwidth} p{0.14\textwidth}}
\toprule
\textbf{Story ID} & \textbf{Raw User Story (Text SRS)} & \textbf{Pre-Prototype Linguistic Defects} & \textbf{Post-Prototype Wireframe Alignment} & \textbf{Resolution Outcome} \\
\midrule
\textbf{US-01} & As a customer, I want to search for suitable storage units quickly. & Vagueness: ``suitable'' and ``quickly'' lack verifiable thresholds. Scope: location and unit size filters are unstated. & Search bar incorporates District/City dropdown, unit size chips ($m^2$), and date picker. Skeleton loader targets $< 2.0$s load time. & \textbf{Resolved} (2/2 defects eliminated) \\
\addlinespace
\textbf{US-02} & As a user, I want to receive notification when my booking is confirmed. & Anaphoric: delivery channel is unstated. Structural: passive voice ``is confirmed'' obscures automated vs. manual approval. & Notification preferences provide checkboxes for Email, SMS, and App Push. Detail screen shows badge: \textit{``Awaiting Manager Approval''}. & \textbf{Resolved} (2/2 defects eliminated) \\
\addlinespace
\textbf{US-03} & As an admin, I want to export all paid and cancelled rental records. & Structural: conjunction ``AND'' leaves set union vs. intersection ambiguous. & Admin dashboard separates actions into \texttt{[Export Paid]} and \texttt{[Export Cancelled]} buttons with date filtering. & \textbf{Resolved} (1/1 defect eliminated) \\
\addlinespace
\textbf{US-04} & As a customer, I want to enter valid card details to pay deposit. & Vagueness: ``valid card details'' lacks field specifications and accepted card brands. & Form specifies Cardholder Name, 16-digit card number with auto-detection, MM/YY expiry, and CVV tooltip. & \textbf{Resolved} (1/1 defect eliminated) \\
\addlinespace
\textbf{US-05} & As a customer, I want to cancel booking in reasonable time and get its deposit back. & Vagueness: ``reasonable time''. Anaphoric: ``its deposit''. Scope: refund percentage is unspecified. & Modal displays tiered policy: $>48$h (100\% refund), $24$--$48$h (50\%), $<24$h (0\%). Calculated refund amount is displayed in VN\DJ. & \textbf{Resolved} (3/3 defects eliminated) \\
\addlinespace
\textbf{US-06} & As a user, the interface must be user-friendly and load storage data efficiently. & Vagueness: ``user-friendly'' and ``efficiently'' are subjective adjectives lacking verifiable criteria. & Design system specifies touch targets $\ge 44 \times 44$px and contrast $\ge 4.5:1$. Static wireframe cannot enforce runtime latency SLA. & \textbf{Partial} (1 latency defect remains) \\
\addlinespace
\textbf{US-07} & As a manager, I want only authorized staff to view facility revenue. & Scope: operational boundary of ``authorized staff'' is indeterminate across team roles. & Role-based matrix restricts revenue view to \texttt{Facility\_Manager} and \texttt{Finance\_Lead}; default denial for general staff. & \textbf{Resolved} (1/1 defect eliminated) \\
\addlinespace
\textbf{US-08} & As a customer, I want the app to show available units every day. & Scope: ``every day'' ambiguously implies either today or interactive calendar dates. & Calendar grid displays monthly availability. Available units show green badges; occupied units are grayed out with end dates. & \textbf{Resolved} (1/1 defect eliminated) \\
\addlinespace
\textbf{US-09} & As a user, I want to save favorite units and share it with partners. & Anaphoric: singular pronoun ``it'' could refer to one unit or the full wishlist collection. & Interface provides two distinct actions: card-level icon \texttt{[Share Unit]} and top-level button \texttt{[Share Full Wishlist URL]}. & \textbf{Resolved} (1/1 defect eliminated) \\
\addlinespace
\textbf{US-10} & As system, automatically send reminder before storage move-in date. & Vagueness: trigger window, delivery hour, and message contents are unspecified. & Automated trigger set to $T-24$h (09:00 AM) containing unit address, access pin code, and digital QR key. & \textbf{Resolved} (1/1 defect eliminated) \\
\bottomrule
\end{tabularx}
\end{table*}

\section{Statistical analysis and Wilcoxon test}

Table~\ref{tab:statistical_summary} summarizes the distribution of ambiguity defects across linguistic categories before and after wireframe alignment.

\begin{table}[htbp]
\centering
\small
\caption{Defect Distribution and Resolution Rates across Categories}
\label{tab:statistical_summary}
\begin{tabular}{l c c c c}
\toprule
\textbf{Ambiguity Category} & \textbf{Pre-Proto ($D_{\text{pre}}$)} & \textbf{Post-Proto ($D_{\text{post}}$)} & \textbf{Eliminated ($\Delta$)} & \textbf{Resolution Rate (ARR)} \\
\midrule
Structural Ambiguity           & 2  & 0 & 2  & \textbf{100.0\%} \\
Scope Ambiguity                & 4  & 0 & 4  & \textbf{100.0\%} \\
Anaphoric Ambiguity            & 3  & 0 & 3  & \textbf{100.0\%} \\
Vagueness and Incompleteness   & 6  & 1 & 5  & \textbf{83.3\%}  \\
\midrule
\textbf{Total}                 & \textbf{15} & \textbf{1} & \textbf{14} & \textbf{93.3\%} \\
\bottomrule
\end{tabular}
\end{table}

Because defect counts represent paired, non-normally distributed ordinal observations, we evaluated the results using the non-parametric Wilcoxon signed-rank test across the ten pairs $(D_{\text{pre}, i}, D_{\text{post}, i})$. All ten pairs exhibited defect reductions ($d_i = D_{\text{pre}, i} - D_{\text{post}, i} > 0$), producing a sum of positive ranks of $W^+ = 55$ and negative ranks of $W^- = 0$. The test statistic yielded $W = 0$, corresponding to a standardized score of $Z = -2.81$ and an asymptotic two-tailed $p$-value of $p = 0.005 < 0.01$. The reduction in ambiguity defects after wireframe inspection is therefore statistically significant.

\section{Discussion and non-functional requirements}

The pilot outcomes reflect several broader patterns observed in requirements engineering literature. Prior to inspection, every user story contained at least one ambiguity defect, matching findings from requirements corpus studies where natural language specifications show defect densities exceeding 70\%. Following wireframe alignment, anaphoric and scope ambiguities were resolved completely. In visual interfaces, form inputs, buttons, and calendar controls occupy distinct spatial positions with unique labels, directly grounding terms that appear ambiguous in prose.

The single unresolved defect occurred in US-06 regarding data loading efficiency. While the wireframe established concrete visual touch targets ($44 \times 44$px) and contrast ratios, static interface mockups cannot represent backend processing time or network latency. This observation aligns directly with the non-functional requirement deficit documented by Rueda et al. (2020) and Abad et al. (2018), where stakeholders focus on visible interface elements while overlooking operational quality attributes.

To resolve this limitation without introducing cumbersome documentation, teams can append explicit Service Level Agreement (SLA) criteria to user story acceptance cards using Gherkin format:
\begin{lstlisting}[basicstyle=\ttfamily\scriptsize, frame=single, backgroundcolor=\color{gray!10}]
Scenario: Facility map loading latency SLA
  Given a customer navigates to the facility floorplan view
  When the client queries storage unit availability from database
  Then all unit status pins must render completely within <= 1.5 seconds.
\end{lstlisting}

\section{Conclusion}

This pilot validation demonstrates that pairing low-fidelity wireframes with structured ambiguity checklists reduces semantic ambiguities by 93.3\% in student project requirements ($p = 0.005$). Grounding textual statements in visual layouts resolves pronoun references and conditional scope effectively. However, because static prototypes do not capture runtime behavior, teams should supplement visual walkthroughs with explicit quality checklists and latency criteria.

\vspace{1.5em}
\hrule
\vspace{1em}
\begin{center}
\footnotesize \textit{Companion technical report for SWR302 (FPT University) --- October 2026.}
\end{center}

\end{document}
